\section{Des collections d'exception au service du droit}
Ancien pôle de référence nationale sous le statut de \enquote{Centre d'acquisition et de diffusion de l'information scientifique et technique} \ac{CADIST}, la bibliothèque Cujas présente une trajectoire singulière \footcite{lobetCujasTraditionModernite2002}. En effet, alors que la majorité des grandes bibliothèques ont constitué leurs fonds patrimoniaux à partir de confiscations révolutionnaires, sa collection d'origine comprend celle \enquote{de la Faculté de Droit de Paris [qui] a été dispersée sous la Révolution} \footcite{balley2ReveillerFonds2012a}. Le fonds actuel a été reconstruit \enquote{rétrospectivement par des bibliothécaires, dans le but de reconstituer \enquote{le cabinet de travail d'un robin d'Ancien Régime} -- telle était, du moins, la doxa transmise jusqu'à présent de bouche à oreille de bibliothécaire} \footcite{balley2ReveillerFonds2012a}. Ce fonds patrimonial s'est essentiellement constitué à partir des années 1870, grâce à des achats, dons ou legs \footcite{balley2ReveillerFonds2012a}. 

Il comprend : \enquote{12 000 volumes imprimés, datant de la fin du XVe siècle au début du XIXe siècle, essentiellement des livres des XVIIe et XVIIIe siècles, 7 incunables, quelques ouvrages précieux postérieurs à 1800 et quelques périodiques. Le quart de ce fonds provient de pays étrangers : Allemagne, Pays-Bas, Angleterre et Italie. Particulièrement riche, il comprend 25 \% d'ouvrages non référencés à la \ac{BNF}. La réserve conserve également des cours et notes de cours manuscrits, ainsi qu'une partie des archives de la Faculté de Droit.} \footcite{balley2ReveillerFonds2012a}. Ces fonds patrimoniaux représentent ainsi les grands courants du droit : droit civil, criminel et canonique, pratique du droit ancien ou nouveau (après 1789), coutumiers, oeuvres des grands jurisconsultes français et étrangers.

La Bibliothèque Cujas regorge d'une richesse documentaire gérée par le \ac{DPCN} représentant \enquote{40 0000 documents patrimoniaux du XVe au XXe siècle conservés dans les magasins patrimoniaux}\footcite{BibliothequeInteruniversitaireCujas}. Elle est forte d'un des plus importants fonds juridiques à l'échelle internationale.

La richesse des collections de Cujas repose historiquement sur des fonds spécialisés et des ouvrages rares. L'établissement est ainsi internationalement reconnu pour conserver un \enquote{ensemble unique de ressources, de collections en sciences juridiques, économiques et politiques}\footcite{BibliothequeInteruniversitaireCujas}. Grâce à cette spécialisation, la bibliothèque est devenue le pôle de référence incontournable pour les chercheurs et les professionnels du droit.

Au-delà de ses monographies, la bibliothèque abrite un précieux gisement de sources primaires comprenant \enquote{les thèses, les mémoires et les publications officielles de la Société des Nations, des Nations Unies, du Journal officiel de la France et des anciennes colonies}\footcite{BibliothequeInteruniversitaireCujas}. À cet inventaire s'ajoute une remarquable \enquote{collection de travaux académiques imprimés en Allemagne [du] XVIe [au] XIXe siècle}, constituve du fonds de Jean-Marie Pardessus. Le fonds Pardessus traite de l'histoire du droit,et notamment du droit germanique. Il contient des thèses allemandes d'institution allemande de l'époque\footcite{BibliothequeInteruniversitaireCujas}. 

Par ailleurs, la Bibliothèque Cujas préserve \enquote{une collection quasi exhaustive de thèses de droit soutenues à la Faculté de droit de Paris}\footcite{BibliothequeInteruniversitaireCujas}. Aujourd'hui encore, l'établissement poursuit activement cette mission patrimoniale de collecte pour les thèses de droit, d'économie et de science politique soutenues au sein des universités Panthéon-Sorbonne (Paris I), et Panthéon-Assas (Paris II)\footcite{HistoireMissionsPartenariats}. Le support de ces travaux varie selon la période, passant de la forme imprimée, à la forme microfichée ou numérique\footcite{BibliothequeInteruniversitaireCujas}.

Les collections de la Bibliothèque Cujas s'étendent bien au-delà des monographies traditionnelles pour intégrer des ressources uniques, telles que les dissertations d'agrégation de droit, les manuels de droit, témoins de la pédagogie du droit, permettant aux chercheurs d'observer et de constater les évolutions du droit français et de son enseignement depuis le début du XIXe siècle \footcite{HoppeDroit}. C'est précisément pour valoriser ce dernier corpus qu'a été conçu le projet \ac{HOPPE-Droit}, financé par le groupement scientifique CollEx-Persée. Grâce à sa richesse patrimoniale, la Bibliothèque Cujas a pu s'appuyer sur cette dynamique collective pour \enquote{sonder l'état de ses collections et enrichir ses données, tout en offrant aux chercheurs un instrument de travail inédit}\footcite{HoppeDroit}. 

Sur le plan méthodologique, \ac{HOPPE-Droit} se définit comme un outil numérique d'analyse relationnelle conçu pour étudier l'évolution de l'enseignement du droit français du XIXe au XXe siècle. La particularité de cet outil réside dans sa capacité à croiser \enquote{des éléments des notices bibliographiques de chacun des ouvrages retenus dans le corpus qu'elle fait résonner les uns par rapport aux autres (titres, auteurs, domaine, éditeurs, collections)}\footcite{HoppeDroit}. À travers la mise en relation par des liens -- tels que les titres, les auteurs, les domaines disciplinaires, les éditeurs et les collections --, l'outil permet de sortir de simples descriptions isolées pour cartographier le squelette conceptuel de la littérature pédagogique de l'Ancien Régime et de l'époque contemporaine \footcite{HoppeDroit}.

Cette modélisation de données massives offre de nouvelles perspectives de recherche, permettant de mettre en évidence les \enquote{tendances historiques qui se sont faites jour en matière d'ouvrages juridiques à vocation pédagogique} et de cartographier \enquote{les réseaux d'édition qui structurent le champ du droit et d'en suivre l'évolution dans le temps}\footcite{HoppeDroit}. Finalement, cette entreprise scientifique n'a pu voir le jour que grâce à une collaboration étroite avec les professionnels des bibliothèques et les chercheurs. Ensemble ils partagent un objectif scientifique commun : \enquote{suivre les évolutions de l'enseignement du droit, saisies à travers le prisme de l'édition juridique à vocation pédagogique} \footcite{HoppeDroit}. 
	
Les fonds spéciaux de la bibliothèque Cujas se caractérisent par leur remarquable diversité et se répartissent en six grands ensembles documentaires distincts. Ces collections comprennent par exemple les publications officielles de la Société des Nations (SDN), celles des Nations Unies (ONU), celles de la collection historique du Journal Officiel (JO), le fonds universitaire de Jean-Marie Pardessus, ainsi qu'un ensemble de cartes postales et les archives historiques de la Faculté de droit de Paris \footcite{BibliothequeInteruniversitaireCujas}. 

Le premier pôle regroupe les sources officielles et étatiques. La collection des publications officielles de la SDN a été constituée dès les années 1920, couvrant l'intégralité de sa production éditoriale entre 1919 et 1946 à travers 534 volumes rédigés en français et en anglais \footcite{BibliothequeInteruniversitaireCujas}. Pour l'ONU, la Bibliothèque Cujas bénéficie du statut de dépositaire depuis 1949, accumulant un fonds colossal de près de 3 800 cartons, dont l'approvisionnement physique a cessé en 2010 à la suite de la dématérialisation des documents officiels \footcite{BibliothequeInteruniversitaireCujas}. Enfin, la bibliothèque conserve l'intégralité du Journal Officiel sur une période s'étendant de 1870 à 2015. Ses numéros historiques sont consultables sous leur forme imprimée jusqu'en 2015, tandis que leur version numérique est accessible en ligne pour les éditions postérieures à 2004 \footcite{BibliothequeInteruniversitaireCujas}.

Le fonds Jean-Marie Pardessus (1772-1853) représente un témoignage capital pour l’histoire universitaire européenne. Ce fonds ancien est entré dans les collections de la Bibliothèque Cujas par l'intermédiaire d'un don consenti par M. de Rozière, petit-fils de Jean-Marie Pardessus, à la Faculté de Droit de Paris \footcite{BibliothequeInteruniversitaireCujas}. 

Enfin, la mémoire de la Faculté de Droit de Paris fait l'objet d'une double valorisation, à la fois iconographique et archivistique. D'une part, la bibliothèque conserve un ensemble original de cartes postales \enquote{caricaturales représentant des professeurs de la Faculté de droit de Paris au début du XXe siècle}, couvrant la période de 1900 à 1914\footcite{BibliothequeInteruniversitaireCujas}, d'autre part, le fonds des archives de la Faculté de Droit de Paris retrace \enquote{la vie institutionnelle, pédagogique et politique de la Faculté de Droit (XIXe-XXe siècle)}\footcite{BibliothequeInteruniversitaireCujas}. Composé de sources administratives, de papiers personnels de professeurs et de manuscrits, ce derniers fonds offre un témoignage précieux sur l'évolution de la vie et de l'organisation universitaire par un tel établissement devenu une bibliothèque d'excellence et de référence mondiale dans le domaine juridique. L'exhaustivité de ses fonds documentaires ne peut se cantonner à un rôle de conservation. C'est précisément cette richesse et cette rareté qui en font une institution majeure au sein du paysage universitaire et qui justifie son intégration dans les réseaux de recherche.